% ID: SEL-D0030-OA-E02
% ESTADO: cotejada. Segunda pasada visual completa y borrador compilado revisados.
% Fuente: Selectividad/matematicas_aplicadas/titular_junio/Examen-Matemáticas Aplicadas a las Ciencias Sociales.pdf
% SHA-256: 5bd7838829b9f653746319907f86642017d12e66ed097db0126fb544a2d0f632
% Página: 1; opción A; ejercicio original 2.

La función de costes de una empresa se puede determinar mediante la expresión
\[
f(x)=40-6x+x^2,\quad\text{para }x\geq0
\]
donde $x$ representa la cantidad producida de un determinado artículo.
\begin{enumerate}[label=\alph*)]
\item \textbf{(1 punto)} \textquestiondown{}Disminuye el coste alguna vez? Determine la cantidad producida de dicho artículo cuando el coste es mínimo y cuál es dicho coste.
\item \textbf{(0.8 puntos)} \textquestiondown{}Cuál sería el coste si no se produjese nada de ese artículo? Si el coste fuese 80, \textquestiondown{}cuántas serían las unidades producidas?
\item \textbf{(0.7 puntos)} Represente gráficamente la función.
\end{enumerate}
